\subsection{本质上积分的定义}

定义 1. --- 对每个函数 \(f \in \mathcal{F}_{+}(\mathrm{T})\) ，把 \(f\) 关于 \(\mu\) 的本质上积分（记作 \(\mu^{\bullet}(f)\) ）定义为数 \(\mu^{*}(f\varphi_{\mathrm{K}})\) （其中 \(\mathrm{K}\) 取遍 \(\mathrm{T}\) 的一切紧子集）所成之集的上确界（有限或否）。设 \(\mathrm{A}\) 是 \(\mathrm{T}\) 的任意子集，令 \(\mu^{\bullet}(\mathrm{A}) = \mu^{\bullet}(\varphi_{\mathrm{A}})\) 。

也使用记号 \(\int^{\bullet}f d\mu ,\int^{\bullet}f(t)d\mu (t),\int^{\bullet}f\mu\) 。

由于 \(f\varphi_{\mathrm{K}} \leqslant f\) 对每个紧子集 \(\mathrm{K}\) （\(\mathrm{T}\) 的）成立，故有

\[
\int^ {\bullet} f d \mu \leqslant \int^ {*} f d \mu .\tag{1}
\]

可能有 \(\mu^{\bullet}(f) \neq \mu^{*}(f)\) ；因为条件 \(\mu^{*}(f) = 0\) 表示 \(f\) 可忽略，而条件 \(\mu^{\bullet}(f) = 0\) 表示 \(f\) 局部可忽略（第四章，§5 第 2 号命题 5），且可能存在并非可忽略的局部可忽略集（第四章，§1 习题 5）。

映射 \(\mu^{\bullet}\) （\(\mathcal{F}_{+}(\mathrm{T})\) 到 \(\overline{\mathbf{R}}\) 中的映射）与 \(\mu\) 在 \(\mathcal{H}_{+}(\mathrm{T})\) 上一致。由此推出，两个测度 \(\mu_{1}\) 与 \(\mu_{2}\) （若 \(\mu_{1}^{\bullet} = \mu_{2}^{\bullet}\) ）相等。

命题 1. --- a) 若 \(f\) 与 \(g\) 是两个局部几乎处处相等的 \(\geqslant 0\) 的数值函数，则 \(\mu^{\bullet}(f) = \mu^{\bullet}(g)\) 。

\begin{enumerate}
\def\labelenumi{\alph{enumi})}
\setcounter{enumi}{1}
\item
  若 \(f\) 与 \(g\) 是两个 \(\geqslant 0\) 的数值函数且 \(f \leqslant g\) ，则 \(\mu^{\bullet}(f) \leqslant \mu^{\bullet}(g)\) 。
\item
  若 \(f\) 是 \(\geqslant 0\) 的数值函数，\(\alpha\) 是 \(\geqslant 0\) 的数，则 \(\mu^{\bullet}(\alpha f) = \alpha \mu^{\bullet}(f)\) 。
\item
  若 \(f\) 与 \(g\) 是两个 \(\geqslant 0\) 的数值函数，则 \(\mu^{\bullet}(f + g) \leqslant \mu^{\bullet}(f) + \mu^{\bullet}(g)\) 。
\item
  若 \((f_n)_{n\in \mathbf{N}}\) 是 \(\geqslant 0\) 的数值函数的递增序列，且 \(f = \lim_{n\to \infty}f_n\) ，则 \(\mu^{\bullet}(f) = \lim_{n\to \infty}\mu^{\bullet}(f_n)\) 。
\end{enumerate}

性质 a)、b)、c)、d) 可由上积分的相应性质立即推出：a) 由第四章 §2 第 3 号命题 6 与第四章 §5 第 2 号命题 5；b)、c)、d) 由第四章 §1 第 3 号命题 10、11、12。为建立 e)，用 K 表示 T 的紧子集所成之集；由第四章 §1 第 3 号定理 3，我们有

\[
\begin{array}{c} \lim _ {n \to \infty} \mu^ {\bullet} (f _ {n}) = \sup _ {n \in \mathbf {N}} \sup _ {K \in \mathfrak {K}} \mu^ {*} (f _ {n} \varphi_ {K}) = \sup _ {K \in \mathfrak {K}} \sup _ {n \in \mathbf {N}} \mu^ {*} (f _ {n} \varphi_ {K}) \\ = \sup _ {K \in \mathfrak {K}} \mu^ {*} (f \varphi_ {K}) = \mu^ {\bullet} (f)  . \end{array}
\]

若 f 与 g 可测，则关系 d) 中的等式成立（第四章 §5 第 6 号定理 5 的推论 4）。更一般地，我们有下述结果：

命题 2. --- 设 \(f,g,h\) 是 \(\mathcal{F}_{+}\) 的三个元素；若 \(g\) 与 \(h\) 可测，则

\[
\int^ {\bullet} f (g + h) d \mu = \int^ {\bullet} f g d \mu + \int^ {\bullet} f h d \mu .\tag{2}
\]

立即可化归为上积分的类似公式的证明。由于 \(f(g + h) = fg + fh\) （按约定 \(0 \cdot (+\infty) = 0\) ），我们有

\[
\int^ {*} f (g + h) d \mu \leqslant \int^ {*} f g d \mu + \int^ {*} f h d \mu ;
\]

剩下的工作是建立反向不等式。设 \(u\) 是使 \(u \geqslant f(g + h)\) 的下半连续函数。令 \(v = \frac{u}{g + h}\) 于使 \(g + h > 0\) 的集上，令 \(v = +\infty\) 于使 \(g + h = 0\) 的集上；则 \(v \geqslant f\) 且 \(u \geqslant v(g + h)\) ，从而

\[
\int^ {*} v (g + h) d \mu \leqslant \int^ {*} u d \mu
\]

并且，由于 v 可测（第四章，§5 第 6 号定理 5 推论 4），

\[
\begin{array}{r l} \int^ {*} f g d \mu + \int^ {*} f h d \mu & \leqslant \int^ {*} v g d \mu + \int^ {*} v h d \mu \\ & = \int^ {*} v (g + h) d \mu \leqslant \int^ {*} u d \mu , \end{array}
\]

由于 u 是任意的，这就蕴含所要的不等式。

推论. --- 设 \(f\) 是 \(\geqslant 0\) 的函数，\((g_n)\) 是 \(\geqslant 0\) 的可测函数的序列；则 \(\int^{\bullet} f\left(\sum_{n} g_n\right) d\mu = \sum_{n} \left(\int^{\bullet} fg_n d\mu\right)\) 。

对有限序列的情形，这是命题 2 的直接推论。无穷序列的情形可借助命题 1 e) 由此推出。

命题 3. --- 对每个有限数 \(\alpha \geqslant 0\) 及 T 上每对测度 \(\mu, \nu\) ，

\[
(\alpha \mu) ^ {\bullet} = \alpha \mu^ {\bullet}
\]

\[
(\mu + \nu) ^ {\bullet} = \mu^ {\bullet} + \nu^ {\bullet}.
\]

此外，关系 \(\mu \leqslant \nu\) 蕴含 \(\mu^{\bullet} \leqslant \nu^{\bullet}\) 。

证明由第四章中的类似命题（§1 第 3 号命题 15）立即得到。

命题 4. --- 对每个在 T 上下半连续的数值函数 \(f \geqslant 0\) ，有 \(\mu^{\bullet}(f) = \mu^{*}(f)\) 。

设 g 是 \(\mathcal{K}_{+}(\mathrm{T})\) 中使 \(g \leqslant f\) 的函数。若 K 是 g 的（紧）支集，则 \(\mu(g) \leqslant \mu^{*}(f\varphi_{\mathrm{K}}) \leqslant \mu^{\bullet}(f)\) 。于是由上积分的定义，\(\mu^{*}(f) \leqslant \mu^{\bullet}(f)\) ，故 \(\mu^{*}(f) = \mu^{\bullet}(f)\) （公式 (1)）。

